\documentclass[11pt,a4paper]{article}
\usepackage[margin=2.4cm]{geometry}
\usepackage{amsmath,amssymb}
\setlength{\parskip}{4pt}\setlength{\parindent}{0pt}
\title{Can multipole cancellation raise the $Q$ of the $\pi$ shift sidewall Bragg cavity?\\[2pt]
\large Honest review and update of the radiation cancellation model (v6.2)}
\author{For Evyatar Rubin, Rosenthal group, Technion}
\date{5 October 2026. Analysis of stored data plus four short verification runs (IGUM 100029, 100034).}
\begin{document}
\maketitle

Labels: \textbf{M} measured from a stored file, \textbf{D} derived from measured values this session,
\textbf{E} expected (theory or estimate). Data: the three $N=98$ far fields of the 20\,$\mu$m study
(original round: \texttt{results\_from\_athena/} \texttt{farfield\_sph\_20um/}; corrected re-run:
\texttt{results\_from\_igum/} \texttt{farfield\_sph\_20um/}); A = TE plain, B = TE overshoot apodization,
C = TM plain. Two independent analyses were made, by Claude and by GPT (Astra), and GPT reviewed the fix; where
they differ it is said. Numbers in items 1 to 5 are from the corrected far field.

\section{Bottom line first}
\begin{enumerate}
\item \textbf{``Cancel one spherical harmonic'' is not the tool for this device.} In A the largest harmonic
group holds 19\,\% of the leak (M), so a perfect cancellation gives at most $1/(1-0.19)=1.24\times$ in
$Q_{\rm rad}$ (D). The five largest groups together hold 68\,\% ($3.1\times$ as a bookkeeping ceiling); in B
the five largest hold 19\,\% ($1.2\times$).
\item \textbf{The harmonic table is not a physical invariant here.} Moving the expansion origin by
0.19\,$\mu$m along $x$ (less than half a pitch) changes A's magnetic dipole from 3\,\% to 19\,\% and $l=3$ from
41\,\% to 23\,\% (D). The grating has no mirror symmetry in $x$ (both arms are built narrow, wide in the
same order), and the standing wave has a node at $x=-0.125$\,$\mu$m and an antinode at $+0.125$\,$\mu$m (M).
There is no natural centre at the scale $1/k=0.17$\,$\mu$m, so ``which harmonic dominates'' depends on a choice.
\item \textbf{The right representation is not a coordinate system. It is the response matrix.} No fixed basis
(spherical, cylindrical, $k_x$) concentrates the leak. The basis that does is device specific: the leak pattern
$\mathbf A_c$ itself is one channel holding 100\,\%. The real question is how much of $-\mathbf A_c$ a
realisable geometry change can produce, Eq.~(\ref{eq:eta}). Both analyses reach this independently.
\item \textbf{Encouraging for the plain TE device.} In A, 83\,\% of the complex leak pattern is reproduced by
12 source patterns located within $\pm0.5$\,$\mu$m of the $\pi$ shift (one tooth each side); extending
the source region to $\pm10$\,$\mu$m adds only 5 points (D, frozen field). A phase scrambled null gives 0.00
and device B gives 0.19 with the same knobs. GPT's falsification test (fit on one part of the sphere, predict
another with no refit) passes two ways and fails one: top cone to side cone $+0.75$, forward half to backward
half $+0.66$, side cone to top cone $-0.41$; for B all are about zero or negative, as they should be. So
``A leaks as \emph{one compact junction source}'' is supported, not proven. It is a synthetic current fit: a ceiling for what junction redesign could remove, not a
prediction of what real teeth achieve (real wall moves couple the current components and are bounded in size).
\item \textbf{The apodized device B is a distributed radiator} (19\,\% reachable within $\pm0.5$\,$\mu$m,
74\,\% within $\pm10$\,$\mu$m). That is a many parameter problem, which is the adjoint inverse design, not a few
knob cancellation. \textbf{Plain TM} is neither: 42 to 50\,\% reachable, two thirds of its power in a grazing
needle ($|u_x|>0.9$), which calls for a periodic feature (the comb).
\item \textbf{The envelope route remains the large lever.} B against A is about $15\times$ in $Q_{\rm rad}$ on
the port numbers, and more if item 8 holds. Mode width alone sets no lower bound on the leak
(Sec.~\ref{sec:env}).
\item \textbf{The far field error was found, explained, fixed and re-measured} (Sec.~\ref{sec:data}). For TE
the corrected far field changes almost nothing: the conclusions above were reached on the faulty data and
survive on the certified data. For TM the low orders reshuffle and the grazing part is still uncertain.
\item \textbf{New finding from the fix: part of the ``loss'' is not cavity radiation.} The power crossing the
surface around the cavity is 7.3\,\% (A), 0.69\,\% (B), 6.5\,\% (C), against port losses $1-T-R$ of 8.55, 2.63
and 8.2\,\% (M). In all three, 1.2 to 1.9 points leave elsewhere, most likely at the two grating ends, which lie
outside the surface (in the short smoke device, with the ends inside, the budget closed to 2\,\% of the loss).
Two consequences, both E until one run with the ends inside confirms it: the apodized device radiates four
times less than its port loss says, and the slow rise of $Q_i$ with length seen in earlier studies is this
constant end loss, not physics of the cavity (Sec.~\ref{sec:q3db}).
\end{enumerate}

\section{The far field error and its correction}\label{sec:data}
On the stitching seam (45$^\circ$ between the monitors) the two projections of the same field give (D, found by
GPT, reproduced by Claude):
\begin{center}\begin{tabular}{lccc}\hline
 & A & B & C\\\hline
complex correlation top/side & 0.85 & 0.80 & 0.41\\
power ratio side/top & 1.31 & 1.32 & 2.26\\\hline
\end{tabular}\end{center}
The larger box (A at 12/12.8) does not repair it (0.88, ratio 1.96).

\textbf{Cause, now proven without any simulation (D).} A flat monitor projection assumes the monitor catches all
the light of its half space. Ours sit 2.15\,$\mu$m from the axis with a $\pm3.4$\,$\mu$m aperture. Pushing the
exact field of a dipole through one such monitor gives the right far field only within about $40^\circ$ of the
monitor normal: 0.73 of the true power at 40 to $50^\circ$, 0.38 at 50 to $60^\circ$, with $\pm15\,\%$ ripple
from $20^\circ$. The side/top ratio of that test follows the measured one bin by bin, and the larger box is
worse in both (monitor at 5.15\,$\mu$m, aperture $\pm6$). So the stitched sphere is good near the two normals
and wrong by tens of percent between them. The handoff statement ``TE box converged'' holds for the stitched
numbers, not for their correctness. TM is the least trustworthy.

\textbf{Fix (implemented).} The two monitors cross. The parts inside the crossing and their mirror images form
one closed tube around the waveguide; $\mathbf E$ and $\mathbf H$ on it at the resonance give the far field in
one coherent surface integral (\texttt{python\_tools/farfield\_surface.py}), with a cosine taper at the open
$x$ ends so the guided wave stays outside the light cone. On analytic dipoles with our exact monitor geometry
it reproduces the exact far field to 0.2\,\% and the absolute radiated power to 0.1\,\%; a deliberately wrong
sign fails. The engine only gains an inert flag that stores the monitor fields
(\texttt{FarFieldConfig.save\_surface\_eh}); geometry and monitors are unchanged. Every result carries its own
certificate: far field power against flux through the surface. GPT reviewed the code, confirmed the formula
and the self test, and found one real bug (a thin strip left open at each corner of the surface, 1.6\,\% in the
field), now fixed and covered by the self test; it estimates the guided wave leakage through the end taper at
$5\times10^{-5}$.

\textbf{Re-run of the three devices with the fix (M, IGUM 100034; same numerics as the original round).}
\begin{center}\small\begin{tabular}{lccc}\hline
 & A (TE plain) & B (TE oversh.) & C (TM plain)\\\hline
$T$, $\lambda$, linewidth, width vs original round & identical & identical & identical\\
far field power / flux through surface & 1.000 & 0.986 & 0.974\\
power through the surface (\% of input) & 7.33 & 0.69 & 6.5\\
port loss $1-T-R$ (\% of input) & 8.55 & 2.63 & 8.21\\
corrected vs stitched far field (correlation) & 0.978 & 0.970 & 0.932\\
largest harmonic, corrected (stitched) & M(1,0) 13.6 (13.2) & M(1,0) 3.6 (3.2) & E(3,0) 8.6 (7.7)\\
share in $l\le5$, corrected (stitched) & 95.5 (93.1) & 30.9 (30.0) & 74.0 (80.1)\\
share at $|u_x|>0.9$ & 9.3 & 10.7 & 66.9\\
standing wave purity $|h|$ & 0.77 & 0.60 & 0.08\\\hline
\end{tabular}\end{center}
Reading. \emph{TE}: the stitched sphere was nearly right after all, because each monitor was only used within
$45^\circ$ of its normal and most of the power lies where one of them is good; individual harmonics move by less
than one point, and the spurious high order tail from the seam disappears. A correlation of 0.98 still leaves
about 20\,\% difference in the field, which matters for cancellation phases, so the corrected field is the one
to use. \emph{TM}: the order of the top harmonics changes, M(2,$\pm$2) drops from 18 to 14\,\%, and the result
depends on the end taper (far field loss 6.6 to 5.9\,\% for tapers of 5 to 25\,$\mu$m) because two thirds of
the power travels within $25^\circ$ of the axis and 10\,\% heads for the open ends. TM needs longer monitors
before its grazing part can be quoted. The handoff's inverse transform (``98\,\% from $\pm5$\,$\mu$m'') is suggestive only; the
compact source fit of item 4 is an independent, per monitor, confirmation of the same picture for A.

\section{Why Johnson's mechanism does not transfer directly}
Johnson's cavity is a point defect with $kR\sim1$: one multipole carries nearly all the loss, the cavity has a
centre, and one parameter moves that one moment through zero. Here:
\begin{itemize}
\item the loss is spread: $\langle l\rangle\approx3.5$ in A, largest group 19\,\%; $\langle l\rangle\approx13$ to 16 in B;
\item a feature at distance $x_p$ from the origin radiates into all orders $l\lesssim kx_p$ (one tooth away is
already $kx_p=1.5$), so no realisable feature addresses a single order;
\item the origin ambiguity of item 2.
\end{itemize}
What does transfer is the mathematics behind it: radiation is linear in the source, so cancellation is a least
squares problem on the whole sphere.

\section{The model, restated around the right objects}
Exact (uniform oxide outside): $\mathbf A(\hat{\mathbf r})=(\mathsf I-\hat{\mathbf r}\hat{\mathbf r})\int\Delta\varepsilon\,\mathbf E\,
e^{-ik\hat{\mathbf r}\cdot\mathbf r}\,d^3r$. For real geometry knobs $\mathbf q$ with response columns
$\mathsf B=\partial\mathbf A/\partial\mathbf q$ and the two constraints (width, resonance) removed by a null space
matrix $\mathsf N$:
\begin{equation}\label{eq:eta}
\min_{\mathbf z}\ \|\mathbf A_c+\mathsf{BN}\mathbf z\|^2_\Omega ,\qquad
\eta^2=\frac{\|\mathsf P_{\mathsf{BN}}\mathbf A_c\|^2}{\|\mathbf A_c\|^2}
\end{equation}
$\eta^2$ is the removable fraction: the squared cosine between the leak and what the knobs can reach. This is the
``basis in which one term dominates'': the first basis vector is $\mathbf A_c/\|\mathbf A_c\|$, and the
singular vectors of $\mathsf{BN}$ say which combinations of knobs radiate independently. The objective must be
$P_{\rm rad}/U$ (normalised to stored energy), not the driven far field alone.

\textbf{What a perturbation must look like to move one harmonic.} To first order
\[
\delta a_{lm}=\underbrace{\int\delta\varepsilon\;\mathbf E_0\cdot\mathbf K_{lm}^*\,d^3r}_{\text{frozen field}}
+\underbrace{\mathcal L_{lm}[\Delta\varepsilon\,\delta\mathbf E]}_{\text{redistribution}},
\]
with $\mathbf K_{lm}\propto j_l(kr)\mathbf X_{lm}$ for a magnetic channel and
$\nabla\times[j_l(kr)\mathbf X_{lm}]$ for an electric one.
Since $j_l\sim(kr)^l$, order $l$ is addressed at $r\sim l/k$: $l=3$ means 0.5\,$\mu$m, the first tooth. The
steepest change of one channel is $\delta\varepsilon\propto-\mathrm{Re}[a_{lm}^*\,\mathbf E_0\cdot\mathbf K^*_{lm}]$,
but it also feeds the neighbouring channels, so selectivity costs one knob per channel touched. At a moved
sidewall use the boundary form ($\mathbf E_\parallel$, $\mathbf D_\perp$), which matters for TE.

\textbf{Counting.} Cancelling $K$ complex channels at fixed width and resonance needs at least $2K+2$ independent
real knobs. The standing wave test $h=\int\mathbf A(\hat{\mathbf r})\cdot\mathbf A(-\hat{\mathbf r})d\Omega/\|\mathbf A\|^2$ gives
$|h|=0.76$ (A), 0.59 (B), 0.20 (C) (D): the leak is not on one phase line, so a single real knob can never cancel
it, and TM least of all.

\section{What each knob can and cannot do (prediction, E unless marked)}
\begin{center}\begin{tabular}{p{3.6cm}p{10.0cm}}\hline
Knob & Assessment\\\hline
Second set of teeth, pure second harmonic $2G$ & Acting on the $\pm\beta$ carriers it creates $\pm3\beta,\pm5\beta$,
all outside the light cone. No direct lever on the leak. It can only act through local (finite support) changes or
through redistribution. Not ``second harmonic cancels the quadrupole''.\\
Innermost one or two teeth each side (width and position), cavity length & The natural knobs for A: the source
is within $\pm0.5$\,$\mu$m. Ceiling 83\,\% (frozen field, D). GPT's independent expectation for the
realisable part is of order 10\,\%; the two are not in conflict (ceiling against realisable) and only a measured
response can decide. Strong resonance and width side effects, to be removed through $\mathsf N$.\\
On axis holes or posts & 2\,\% (Claude) and 0.6 to 4.7\,\% (GPT) of A's leak, D. Dead end, consistent with v6.\\
Tooth position shifts along the arms & Continuous phase, response $\propto ik_x\tilde{\mathbf J}$. The natural
knobs for B's distributed leak. This is what the inverse design already uses.\\
Periodic comb in the cladding & Narrow band in $u_x$ with a phase dial. Suits the TM grazing needle (62\,\% of C's
power at $|u_x|>0.9$, D), not the broad TE hump at $|u_x|\approx0.75$.\\\hline
\end{tabular}\end{center}
Which current component carries A's leak is \emph{not} established: on the full sphere the best single family is
longitudinal sidewall current odd in $y$ (57\,\%), but inside the top monitor cone the transverse families explain
the same field. Only the compactness is robust.

\section{Envelope view and the absence of a width bound}\label{sec:env}
The source carriers sit at $\pm\beta$, just outside the light cone ($\beta-k=0.46$\,$\mu$m$^{-1}$); the leak is
the envelope's Fourier content beyond that offset plus the junction term. For the same 20\,$\mu$m intensity width
the scalar light cone fraction is $9\times10^{-5}$ for $e^{-\gamma|x|}$ and $3\times10^{-15}$ for a Gaussian (D,
GPT); a $\mathrm{sinc}^2$ envelope has none at all. So the width specification costs nothing in principle. The
limits are practical: a Gaussian needs $\kappa=0.069$\,$\mu$m$^{-1}$ at the half width (plain device: about
0.035 to 0.044), the joins, the finite length, and the transverse currents the scalar picture omits. Measured A
does not look like a pure Lorentzian tail (51\,\% of the power at $0.5<|u_x|<0.8$, only 3\,\% above 0.95), which
is again the junction source of item 4. B is not shown to be near any bound.

\section{Corrections to v6.2 and to the handoff}
\begin{itemize}
\item v6.2 Sec.~4.4 (TE: $E_y$ even in $x$, electric dipole allowed, magnetic dipole forbidden, features must sit
at the centre or in $\pm x$ pairs): not valid. There is no $x$ mirror symmetry; both parity sectors hold about half
the power (D, GPT). Claude's first hypothesis ($E_y$ odd, central node) was also wrong for the same reason.
\item v6.2 ceilings of 10 to 14\,\%: reproduced, but they are the ceiling of \emph{on axis} $y$ dipoles only,
not of the junction teeth.
\item Handoff ``only dipole is M(1,0), a pillar cannot cancel it'': origin dependent and too categorical; the CSV
also has 0.6\,\% electric dipole.
\item Handoff ``+0.011 in $T$ for cancelling M(1,0)'': arithmetic confirmed (+0.0109, GPT), assumptions as above.
\item Formulas (1) to (11) of v6.2: no error found by either analysis.
\end{itemize}

\section{Plan}
\textbf{Done on 5 October.} Far field error diagnosed and fixed; three devices re-measured with the corrected
projection (IGUM 100034); GPT review of the fix applied.

\textbf{Next, in this order (nothing of this is dispatched).}
\begin{enumerate}
\item \textbf{One run, ends inside:} plain TE at $N=98$ with the far field monitors spanning the whole device
(110\,$\mu$m instead of 80). Decides item 8: if the surface then carries the full 8.55\,\%, the missing 1.2 points
are end scattering. It also gives TM style grazing light a longer path, so the same setting is the one to use
for TM.
\item \textbf{One run, junction source:} the same device with the complex vector field stored in a small volume
$\pm2$\,$\mu$m around the $\pi$ shift, both sides of every sidewall. This measures the junction source directly
and settles which current component radiates. Can be the same run as item 1.
\item \textbf{One run, one perturbation:} the single combined width and position change of the innermost teeth
that the constrained model of Eq.~(\ref{eq:eta}) selects. Compare predicted and simulated radiation change,
resonance shift, width change and $P_{\rm rad}/U$. Agreement calibrates that direction including
redistribution; disagreement kills the frozen field ceiling of item 4.
\item \textbf{Only if item 3 agrees:} four to six junction knobs, $\pm h$ each (about ten short runs), solve
Eq.~(\ref{eq:eta}), one unseen validation geometry, then one long device to confirm the Q3dB transfer of
Sec.~\ref{sec:q3db}. Success criterion fixed in advance.
\item For B, no few knob study: give the inverse design the radiated power objective, computed with the
corrected projection.
\end{enumerate}
Known loose end in the code: the extraction pulls the whole recorded band into Python to read the wavelengths
(the engine's existing pattern; ran at 120\,GB). Slice it inside Lumerical before using more frequency points.

\section{What this means for the Q3dB devices (prediction)}\label{sec:q3db}
The Q3dB devices are longer ($N\approx166$ to 172 against 98 here). The transfer rests on two facts.

\textbf{1. At the $-3$\,dB point the loaded $Q$ is a fixed fraction of the radiation $Q$.} With
$T=(Q_L/Q_c)^2$ and $1/Q_L=1/Q_c+1/Q_i$, $T=\tfrac12$ gives
\[
Q_{3\rm dB}=\bigl(1-\tfrac1{\sqrt2}\bigr)\,Q_i=0.2929\,Q_i ,
\]
which reproduces the four measured anchors to about 2\,\% (earlier study). Length only sets the coupling $Q_c$;
it is re-tuned to stay at $-3$\,dB.

\textbf{2. The leak of the plain device is a junction source} (item 4): it does not know how long the mirrors
are, so $Q_i$ should not depend on $N$. The port numbers seemed to contradict this: $Q_i=Q_L/(1-\sqrt T)$ gives
$3.8\times10^4$ at $N=98$ against $4.4\times10^4$ at the TE Q3dB point ($N=166$, $Q_L=12\,903$). Item 8 resolves
it: at $N=98$ only 7.33 of the 8.55 points of loss are cavity radiation. With the cavity part alone,
$Q_i=4.35\times10^4$ at $N=98$ (D), equal to the $N=166$ value ($4.41\times10^4$) within about 1\,\%. The drift was the constant end
loss, which weighs 14\,\% of the loss at $T=0.91$ and almost nothing at $T=0.5$. So a response measured on the
cheap $N=98$ device should carry over, and the 0.2929 rule is safe at the $-3$\,dB point. For the overshoot
device the same correction gives $Q_i\approx2\times10^6$ instead of $6\times10^5$ (E, rests on item 8): its
Q3dB version would be correspondingly better than the port based extrapolation, and would run into the
measurement limit above $Q\sim5\times10^4$. One run of a long plain device with the surface fields stored
checks both statements at once.

\subsection*{What happens to the harmonic percentages when the device is lengthened}
\textbf{Prediction (E, not yet measured): for the plain devices the percentages stay the same.} The percentages
describe the \emph{shape} of the radiation pattern, and the shape is set by where the light leaves and with
what field. For the plain TE device that is the junction, within $\pm0.5$\,$\mu$m of the $\pi$ shift. Adding
mirror periods at the outer ends changes how strongly the cavity is coupled to the bus, not the field at the
junction, so the same pattern is radiated, only stronger. The TE Q3dB device is exactly device A lengthened
from 98 to 166 periods (same pitch, corrugation and width).
\begin{center}\small\begin{tabular}{p{5.3cm}p{9.2cm}}\hline
Quantity & From $N=98$ to the Q3dB length\\\hline
Harmonic percentages, plain TE & unchanged within about 1 to 2 points (E). Table of Sec.~\ref{sec:data}: M(1,0) 13.6,
E(3,$\pm$3) 19.0, M(2,$\pm$1) 12.6, E(4,$\pm$4) 12.2, E(2,$\pm$2) 11; $l\le5$ holds 95.5\,\%\\
Radiation $Q_i$ & unchanged, $4.4\times10^4$ (D, item 8)\\
Radiated share of the input power & rises from 7.3\,\% to about 41\,\% at $-3$\,dB ($1-T-R$ with $T=0.5$,
$R=(1-\sqrt T)^2$): same pattern, more of it\\
Standing wave purity $|h|$ & may rise from 0.77: a more weakly coupled cavity is closer to a pure standing wave.
This is the one place a small reshuffle could come from\\
End scattering & not part of the cavity pattern at any length; negligible share of the loss at $-3$\,dB\\
Overshoot device & pattern unchanged \emph{if} the extra periods are plain mirror periods outside the 60
apodized ones: its sources sit at $\pm6$, $\pm13$, $\pm25$\,$\mu$m, all inside the apodized part (E)\\
Plain TM & expected unchanged for the same reason, but its grazing needle is not yet certified, so this is the
weakest of the three statements\\\hline
\end{tabular}\end{center}
What would change the percentages is a change of the junction or of the envelope (corrugation, apodization,
any cancelling feature), not length. This matters for the plan: if it holds, the harmonic content and every
cancellation response can be measured on the short, cheap device and applied to the Q3dB one. \textbf{It is a
prediction and needs one run to confirm:} plain TE at $N=166$ with the surface fields stored, compared harmonic
by harmonic with the $N=98$ table above. Pass criterion, fixed now: every harmonic group within 2 points and
far field correlation above 0.95.

\subsection*{What a cancellation is worth at the $-3$\,dB point}
If a junction change removes a fraction $f$ of the leak at fixed width, then $Q_i\to Q_i/(1-f)$ and the device
must be lengthened by $\Delta N=\ln[1/(1-f)]/(2\kappa\Lambda)$ periods per side to return to $-3$\,dB
($\kappa=0.034$\,$\mu$m$^{-1}$, $\Lambda=0.5$\,$\mu$m). For the TE Q3dB device ($Q=12\,903$ at $N=166$):
\begin{center}\begin{tabular}{lcccc}\hline
removed fraction $f$ & 0.10 & 0.30 & 0.50 & 0.80\\\hline
predicted $Q_{3\rm dB}$ (E) & 14\,300 & 18\,400 & 25\,800 & 64\,500\\
extra periods per side (E) & 3 & 10 & 20 & 47\\\hline
\end{tabular}\end{center}
The mode width is set by $\kappa$ and stays at 20\,$\mu$m to first order; the junction change itself must be
width neutral, which is what the null space constraint in Eq.~(\ref{eq:eta}) enforces. The existing q3db
predictive engine then gives the new length in one confirmation run. For comparison, the comb pair device
reached $Q=18\,093$ against 13\,930, equivalent to $f\approx0.23$ on TM. Above $Q\sim5\times10^4$ the measurement itself needs
the finer grid and longer ring down (known trap), so the $f=0.8$ column could not be verified with the standard
numerics.

\section{Is there a basis with one dominant term? (asked of both models)}
Both answers agree.
\begin{itemize}
\item A change of basis preserves inner products. It can always put the whole leak into one coefficient
($\mathbf e_1=\mathbf A_c/\|\mathbf A_c\|$), but that is a tautology: it cannot increase the overlap between the
leak and what a physical knob radiates.
\item The meaningful ``one term'' is a \emph{controllable} pattern. Take the singular value decomposition of the
constrained, real valued response matrix, $\mathsf{BN}=\mathsf U\Sigma\mathsf V^{\!\top}$, with the columns scaled
by the allowed geometry change. If one left singular vector with a large singular value captures most of
$\mathbf A_c$, the device has a single controllable channel, and the matching right singular vector \emph{is} the
perturbation: the list of tooth changes to apply. This is the generalisation of Johnson's idea that fits this device.
\item To isolate the cavity's own pattern from background scattering, fit the complex far field across the
resonance to a pole, $\mathbf A(\omega)=\mathbf d/(\omega-\tilde\omega)+\mathbf A_{\rm bg}$; the residue
$\mathbf d$ is the pattern to cancel. The far field monitors already record a band (81 points); only the
projection at more than one wavelength has to be saved.
\item Cylindrical harmonics about $x$ or a $k_x$ spectrum are convenient for the long apodized device, but they
do not concentrate the leak either.
\end{itemize}
Where the two analyses differ: Claude's frozen field ceiling for the junction of A is 83\,\%; GPT expects
the realisable part to be of order 10\,\% and regards the ceiling as optimistic until a real wall move is
measured. Items 2 and 3 of the plan decide between them.

\appendix
\section{The analytical model of v6.2, carried over (with corrections marked)}
This appendix keeps the content of the earlier note (``Analytical model for radiation channel cancellation'',
v6.2, 27 September 2026) so that this document stands alone. Derivations are in v6.2, Appendix A; neither
analysis found an error in the formulas.

\textbf{Setting.} Axis $x$, transverse $y$, vertical $z$, time dependence $e^{-i\omega t}$. $n_{\rm SiN}=1.97$,
$n_c=1.444$; uniform oxide everywhere else. At 1559\,nm: $k=n_ck_0=5.82$\,$\mu$m$^{-1}$, $1/k=0.17$\,$\mu$m;
$\beta=6.28$ (TE, pitch 500) and 6.08\,$\mu$m$^{-1}$ (TM, pitch 517), just outside the light cone. Total contrast
$\Delta\varepsilon=\varepsilon-n_c^2=1.80$ in SiN; a feature changes it by $\delta\varepsilon_P=\mp1.80$ (oxide
hole in SiN, SiN post in oxide).

\textbf{Exact far field.} Moving the contrast to the right of Maxwell's equations makes the core a source in
uniform oxide:
\[
\mathbf E\to\frac{k_0^2e^{ikr}}{4\pi r}\mathbf A(\hat{\mathbf r}),\quad
\mathbf A=\mathcal R[\Delta\varepsilon\mathbf E]=(\mathsf I-\hat{\mathbf r}\hat{\mathbf r})\!\int\!\Delta\varepsilon\,\mathbf E\,
e^{-ik\hat{\mathbf r}\cdot\mathbf r}d^3r,\quad
\frac{dP}{d\Omega}=\frac{k_0^4}{32\pi^2\eta}|\mathbf A|^2,\quad Q_{\rm rad}=\frac{\omega U}{P}.
\]
$\mathbf A$ is the Fourier transform of the source on the sphere $|\mathbf q|=k$. $\mathcal R$ is linear;
everything else is bookkeeping.

\textbf{Multipole form.} With $\mathbf X_{lm}=\mathbf LY_{lm}/\sqrt{l(l+1)}$,
$\mathbf A=\sum[\alpha_{lm}\mathbf X_{lm}+\beta_{lm}\hat{\mathbf r}\times\mathbf X_{lm}]$ and
$P\propto\sum(|\alpha_{lm}|^2+|\beta_{lm}|^2)$: magnetic and electric multipoles, powers add (Johnson's
argument). Near field form:
\[
\alpha_{lm}=4\pi(-i)^l\!\int\!\Delta\varepsilon\,j_l(kr)\,\mathbf X^*_{lm}\!\cdot\mathbf E\,d^3r,\qquad
\beta_{lm}=\frac{4\pi i(-i)^l}{k}\!\int\!\Delta\varepsilon\,\mathbf E\cdot\nabla\times[j_l(kr)\mathbf X^*_{lm}]\,d^3r .
\]
Since $j_l\sim(kr)^l$, a source of radius $R$ needs $l\sim kR$. Dipole shares without any harmonics:
$f_{\rm ED}=\frac{3}{8\pi}|\int\mathbf A\,d\Omega|^2/\int|\mathbf A|^2d\Omega$, and the same with
$\hat{\mathbf r}\times\mathbf A$ for $f_{\rm MD}$.
\emph{Correction:} v6.2 fixed the origin at the cavity centre and assumed $E_y$ even in $x$. The grating has no
$x$ mirror symmetry and the shares move with the origin at the 0.2\,$\mu$m scale (item 2); only the $y$ and $z$
parities are selection rules (TE: electric $l+m$ even, magnetic $l+m$ odd; TM the converse).

\textbf{Adding a feature: exact first order split.}
\[
\mathbf A_{\rm new}-\mathbf A_c=\underbrace{\mathcal R[\delta\varepsilon_P\mathbf E_0]}_{\text{frozen field}}
+\underbrace{\mathcal R[\Delta\varepsilon\,\delta\mathbf E]}_{\text{redistribution}}
+\delta k\,\partial_k\mathcal R[\Delta\varepsilon\mathbf E_0],\qquad
(\nabla\!\times\!\nabla\!\times-k_0^2\varepsilon)\delta\mathbf E=k_0^2\delta\varepsilon_P\mathbf E_0+2k_0\delta k_0\varepsilon\mathbf E_0 .
\]
The first term is the feature's own radiation in the unchanged cavity field (closed form, used for screening).
The second is the change of the cavity field caused by the feature; it has no closed form and is what a finite
difference run measures. Both are first order.

\textbf{Small feature} ($R\lesssim0.15$\,$\mu$m): a point dipole,
\[
\mathcal R[\delta\varepsilon_P\mathbf E_0]\simeq(\mathsf I-\hat{\mathbf r}\hat{\mathbf r})\frac{\mathbf p}{\varepsilon_0}
e^{-ik\hat{\mathbf r}\cdot\mathbf r_p}S(\hat{\mathbf r}),\qquad
\mathbf p=\varepsilon_0\,\delta\varepsilon_PV_PL\,\mathbf E_0(\mathbf r_p),
\]
with $S$ the shape factor of the feature. Local field factor $L=2\varepsilon_{\rm SiN}/(\varepsilon_{\rm SiN}+\varepsilon_c)=1.30$ for a vertical oxide
cylinder in SiN; at a sidewall use the core side $E_y$ (the cladding side value is $1.86\times$ larger). A
feature at the origin is a pure dipole; displaced by $x_p$ its power spreads over the orders:
\begin{center}\small\begin{tabular}{cccccccc}\hline
$kx_p$ & $l=1$ & 2 & 3 & 4 & 5 & 6 & 7\\\hline
0 & 1.00 & 0 & 0 & 0 & 0 & 0 & 0\\
1 & 0.86 & 0.13 & 0.01 & 0 & 0 & 0 & 0\\
2 & 0.55 & 0.36 & 0.08 & 0.01 & 0 & 0 & 0\\
4 & 0.14 & 0.29 & 0.33 & 0.18 & 0.05 & 0.01 & 0\\
6 & 0.06 & 0.10 & 0.16 & 0.26 & 0.24 & 0.13 & 0.04\\\hline
\end{tabular}\end{center}
One tooth from the centre is already $kx_p\simeq1.5$. This table is why a feature cannot be told to radiate one
order.

\textbf{Pair, comb, phase.} Two equal dipoles at $\pm x_p$:
$\mathbf A_P=2\cos(kx_p\sin\theta\cos\varphi)(\mathsf I-\hat{\mathbf r}\hat{\mathbf r})\mathbf p/\varepsilon_0$; a
sign, no continuous phase. A periodic comb of period $\Lambda_2$ and offset $\delta x$ under an envelope $f$:
\[
\mathbf A_P\propto\Delta\varepsilon_2\,e^{i\phi_2}\,\widetilde{(gf)}\bigl(k\sin\theta\cos\varphi-(\beta-K_2)\bigr),\qquad
\phi_2=\frac{2\pi\delta x}{\Lambda_2},\quad u_x=\frac{\beta-K_2}{k},\quad\Lambda_2=\frac{2\pi}{\beta+k|u_x|}.
\]
It emits a narrow band around $u_x$ with a continuous phase set by its offset. A feature inside the core is
driven by a standing wave: a sign and a size, nothing else, so it moves a channel amplitude along one line of
the complex plane. Only a comb offset or a position in the cladding gives a phase. Moving a boundary is not a
small volume change: use the continuous fields $\mathbf E_\parallel$ and $D_\perp$ (Johnson 2002).

\textbf{Cancellation condition, ceiling, when it hurts.} With feature patterns $\mathbf b_j$ and real strengths
$q_j$, $P\propto\|\mathbf A_c+\sum q_j\mathbf b_j\|^2$. One feature:
\[
q_*=-\frac{\mathrm{Re}\langle\mathbf b,\mathbf A_c\rangle}{\|\mathbf b\|^2},\qquad
\eta^2=\frac{[\mathrm{Re}\langle\mathbf b,\mathbf A_c\rangle]^2}{\|\mathbf b\|^2\|\mathbf A_c\|^2},\qquad
\frac{P(q)}{P_c}=1-\eta^2+\eta^2\Bigl(\frac q{q_*}-1\Bigr)^2 .
\]
Several features: $\mathrm{Re}(\mathsf G)\mathbf q=-\mathrm{Re}\,\mathbf b$ with
$G_{ij}=\langle\mathbf b_i,\mathbf b_j\rangle$. Cancelling a fraction $a$ of the amplitude of a channel holding
a fraction $f$: $P_{\rm new}/P_c=(1-f)+f|1-a|^2$; 90\,\% gives at most $10\times$, 50\,\% gives $2\times$, 8\,\%
gives $1.09\times$. The loss rises above the baseline for $q>2q_*$ or the wrong sign.

\textbf{Mode width (two mechanisms).} A feature shifts the resonance,
\[
\frac{\Delta\omega}{\omega}\simeq-\frac12\,\frac{\delta\varepsilon_PLV_P|\mathbf E_0(\mathbf r_p)|^2}{\int\varepsilon|\mathbf E_0|^2d^3r},
\]
which changes the decay rate $\gamma=\sqrt{\kappa_g^2-\delta_\beta^2}$ (intensity width $\ln2/\gamma$) only at
second order at gap centre. Features repeating with the pitch change $\kappa$ at first order: oxide holes in the
narrow segments shorten the mode, in the wide segments lengthen it; cladding posts at 1.5 to 2\,$\mu$m do neither
to a measurable degree. Constraint and feasibility: with $\mathsf N$ spanning the null space of the width and
resonance rows,
\[
\mathsf C=\begin{bmatrix}\mathrm{Re}(\mathsf{BN})\\\mathrm{Im}(\mathsf{BN})\end{bmatrix},\quad
\mathbf y=\begin{bmatrix}\mathrm{Re}\,\mathbf d_0\\\mathrm{Im}\,\mathbf d_0\end{bmatrix},\quad
\mathbf z_*=-\mathsf C^+\mathbf y,\quad
\rho=\frac{\|(\mathsf I-\mathsf{CC}^+)\mathbf y\|^2}{\|\mathbf y\|^2},
\]
where $\rho$ is the fraction the family cannot remove at fixed width. Cancelling one complex channel under both constraints is four real
conditions.

\textbf{Results of v6.2 that stand (they do not depend on the far field files).}
\begin{itemize}
\item Comb period: the TM comb emission is predicted at $u_x=-0.989$ (531\,nm), $-0.970$ (536), $-0.936$ (545),
$-0.915$ (551); measured: 531 locked ($+16.3\,\%$ $Q$ at fixed width), 536 slightly worse, 545 harmful, 551 null.
\item Phase circle: $T(0^\circ,90^\circ,180^\circ,270^\circ)=0.8694,0.8586,0.8689,0.8797$ fits a pure first
harmonic in the comb offset, as $e^{i\phi_2}$ requires; lock at $270^\circ$.
\item TE post pair near the centre (two $r=80$\,nm SiN posts at $x=0$ and 270\,nm, $y=\pm700$\,nm): loss 0.119 to
0.093 at fixed width, a near uniform scaling of the whole pattern, that is the redistribution term, not the
cancellation of a channel.
\item Numbers: $L=1.30$; $\varepsilon_{\rm SiN}/\varepsilon_c=1.86$; TM evanescent decay 1.75\,$\mu$m$^{-1}$;
superposition holds to 4.4\,\% for $r=80$\,nm posts at 1.08\,$\mu$m spacing; linewidths 1.1\,nm ($N=80$),
0.12\,nm ($N=166$).
\end{itemize}
\textbf{Results of v6.2 that are superseded.} Its channel content for TE $N=80$ (8\,\% dipole, 28\,\% in $l=3$) and
its removable fractions (centre feature 0.08, pairs up to 0.14) were computed on far field files that failed the
consistency test, about an assumed centre, and for on axis $y$ dipoles only. Use Sec.~\ref{sec:data} and item 4
instead. Its remedy for the far field (closed surface with $\mathbf E$ and $\mathbf H$, guided modes handled at
the ends) is what was implemented here.

\section{Remaining content of v6.2}
Everything else that v6.2 contained, so that nothing is lost. Statements that today's work changes are marked.

\textbf{What the model is for.} FDTD answers one geometry in hours and gives no reason. The model does three
things FDTD cannot: it splits the loss into channels whose powers add, so the value of removing a channel is
known before any run; it maps geometry to radiation linearly, so a cancelling feature is found by a small linear
solve, not a sweep; and it gives a smooth, cheap objective for the inverse design in place of a noisy $Q$. The
deliverable is a calibrated local model around the seed, with $\mathbf q$ a short vector of real geometry changes:
\[
\mathbf d\simeq\mathbf d_0+\mathsf B\mathbf q,\qquad L\simeq L_0+\boldsymbol\ell^{\!\top}\mathbf q,\qquad
\omega\simeq\omega_0+\mathbf f^{\!\top}\mathbf q,\qquad
\min\|\mathbf d_0+\mathsf B\mathbf q\|^2\ \text{with}\ \boldsymbol\ell^{\!\top}\mathbf q=0,\ \mathbf f^{\!\top}\mathbf q=0,
\]
where $\mathbf d$ holds the complex channel amplitudes normalised so that $\|\mathbf d\|^2=\omega/Q_i$ and $L$ is
the mode width. This is the structure of the ns2 adjoint step with radiated power in place of transmission.
``Harmonic'' has three meanings: angular multipoles $(l,m)$ label far field patterns, grating harmonics
$2\pi p/\Lambda$ label the geometry, optical harmonics are not part of this work; ``channel'' always means an
angular pattern. Definitions: $Q_L$ from the transmission linewidth; $Q_i$, $Q_c$ from $T$ and $R$ at resonance;
mode width = half maximum above floor of the envelope of the $y$ integrated $|\mathbf E|^2$ at $z=0$.

\textbf{When cancellation hurts (the three reasons).} The feature's pattern is not the channel it targets and
feeds others; the target channel is small, so $q_*$ and the safe window are small; the feature changes the
width, and restoring it through the corrugation rewrites $\mathbf A_c$. For TE the risk is larger than for TM:
$E_y$ peaks on the axis and is enhanced by $L=1.3$ inside a hole, so a centre hole is a strong radiator and a
strong detuner while the dipole channel it addresses is small. Why ``cancel $l=3$'' is no better than ``cancel
the dipole'': to cancel the cavity's $l=3$ with the pairs available, most of the pair's power lands in $l=2$ and
$l=4$, where the cavity also radiates but with unrelated phases (a pair at 0.75\,$\mu$m has 4\,\% of its own
power in $l=3$, 33\,\% in $l=2$, 30\,\% in $l=4$).

\textbf{Mode width, details.} Plain device $\kappa_g\simeq0.044$\,$\mu$m$^{-1}$ (v6.2; 0.034 from the q3db
calibration at corr 250). The resonance shift of a feature is blue for oxide; at exact gap centre it changes the
width only at second order (0.6\,\% of width per nm for the plain device, quadratic), at first order if the seed
sits off centre or the local $\kappa(x)$ near the centre is comparable to the detuning. The 0.365\,$\mu$m/nm in
the adjoint code is a path derivative fitted along an optimiser trajectory, not this derivative; a gap length
sweep (\texttt{cavity\_neg\_detuning\_nm}) measures it in three to five runs. Holes at the tooth transitions add
in quadrature (a smaller change); a detuned period beats and changes the average $\kappa$ little while radiating
at the folded angle. The comb: width 19.97 to 19.90\,$\mu$m, $\lambda$ $+10$\,pm. For a centre feature the cheap
way to hold the resonance is the gap: $L_{\rm gap}=-(f_C/f_G)A_C$,
$\mathbf b_{\rm eff}=\mathbf b_C-(f_C/f_G)\mathbf b_G$,
$A_{C*}=-\mathrm{Re}\langle\mathbf b_{\rm eff},\mathbf A_c\rangle/\|\mathbf b_{\rm eff}\|^2$, and the remaining
width change $(\ell_C-\ell_Gf_C/f_G)A_{C*}$ is checked against the measured envelope, which gates every decision.
A comb supplies two controls (size, offset), a centre hole one. Measured width signs (TM hole lattice against a
15.5\,$\mu$m control): holes in the narrow segments 10.2\,$\mu$m (corrugation trimmed); shifted a quarter pitch
12.5; period 545: 16.0 with $Q$ 1209 to 918; untrimmed matched lattice: resonance destroyed; in core 531\,nm comb
on the 517\,nm pitch: 25.6\,$\mu$m.

\textbf{Procedure of v6.2 for connecting model and simulation} (the Plan section above is its updated form).
(1) Baseline run, port driven, run to full decay; the far field frequency must sit on the resonance (0.07\,nm off
is 6\,\% of the $N=80$ linewidth and half the $N=166$ linewidth). (2) A far field that passes the consistency
test: closed surface with $\mathbf E$ and $\mathbf H$, guided modes handled at the ends, window on the long
faces, transform in Python; validate that far field power equals box flux and that the result is stable when the
ends move. \emph{Done today, with the lateral surface and taper.} (3) Channel analysis: one channel or region
above $f\simeq0.7$ means a few features; one compact region with many $(l,m)$ means an aimed comb; spread means
the model is the optimiser's objective. (4) Screen without runs from the baseline field. (5) Calibrate: gap
sweep first, then $\pm h$ pairs per kept coordinate with the resonance tracked, both mesh levels, the noise
floor rule, the control run with the scatterer explicitly off; later an adjoint with a plane wave source from
the target direction gives whole rows of $\mathsf B$. (6) Solve and validate on one unseen combined geometry;
report predicted and simulated complex amplitudes, $Q_i$, $Q_c$, width and resonance together; a higher peak $T$
alone is not evidence. Three remarks on observables: peak transmission depends on the balance of port coupling
and internal loss, so a higher peak can come from changed coupling; a $Q$ peak with extra nodal lines is
supporting evidence, the targeted coefficient passing through zero with a loss budget is the proof; a fixed
optical mode width keeps the spatial averaging of the ultrasound field fixed, but not the strain optic overlap or
the optical bandwidth, which shrinks as $Q$ rises. Suppressing one channel does not create a bound state in the
continuum.

\textbf{Devices and numbers quoted in v6.2.} Plain TE $N=80$: corr 300, $Q_L=1411$, width 15.6\,$\mu$m,
$T=0.876$. Plain TE $N=166$: corr 250, $Q_L=12\,903$, width 20.5\,$\mu$m. Overshoot seed (Nt60/W20):
$Q_i=6.1\times10^5$ at $N=98$, $1.16\times10^6$ at $N=130$ (port based; see item 8), width 19.7\,$\mu$m, strongly
overcoupled. TM evanescent tail: 3.6\,\% at 1.9\,$\mu$m, 0.5\,\% at 3\,$\mu$m, which is why the 551\,nm comb at
standoff did nothing. The TE comb: 590\,nm aims at $u_x=-0.75$, 580\,nm at 0.78. Born test for superposition:
spacings of 270 and 1080\,nm pass, 135, 405, 540 fail. One seed run costs 0.3 to 0.6 GPU hours, a ten point
calibration 3 to 6.

\textbf{v6.2's own far field test and indicative numbers (superseded by Sec.~\ref{sec:data}).} In the octant
$u_y,u_z>0.2$, $|u_x|<0.9$ its files gave a top/side magnitude correlation of 0.12 (TE) and 0.22 (TM) and power
ratios 0.22 and 0.68; it blamed single plane projections with the bus waveguide leaving through the plane ends.
Its indicative TE $N=80$ content: $f_{\rm ED}=0.08$, $f_{\rm MD}=0.02$; by order $l=1..8$: 0.10, 0.10, 0.28,
0.19, 0.10, 0.08, 0.02, 0.02; 55\,\% of the power at $0.6<|u_x|<0.9$. Its removable fractions (frozen field,
complex knob): centre feature 0.08; pairs at $\pm0.25$ to $\pm1.5$\,$\mu$m 0.005 to 0.06; centre plus two pairs
0.12; centre plus seven pairs to 3\,$\mu$m 0.14 (real knobs 0.10 best, 0.04 worst); two row post comb 0.02. Its
conclusion, ``frozen field cancellation by small features is worth at most about 10\,\% for TE'', holds for the
on axis features it tested and is replaced by item 4 for the junction teeth.

\textbf{Numerical checks of v6.2.} The multipole formulas, the dipole share formulas and the point dipole
formula were checked against an independent computation (four random dipoles, far field on a Gauss Legendre
grid, projection to $l=14$): agreement to $6\times10^{-9}$, Parseval to six digits. Scripts delivered with v6.2:
\texttt{verify\_multipole.py}, \texttt{check\_l1.py}, \texttt{check\_l1b.py}, \texttt{ff\_inspect.py},
\texttt{ff\_channels.py}, \texttt{ff\_features.py}. Not verified there: the local field factor for a real finite
hole.

\textbf{Derivation (v6.2 Appendix A, in outline).} (1) Maxwell with the contrast on the right:
$\nabla\times\nabla\times\mathbf E-k^2\mathbf E=k_0^2\Delta\varepsilon\mathbf E$, $k^2=\varepsilon_ck_0^2$; nothing
approximated. (2) Outgoing solution with the oxide Green tensor
$\mathsf G=(\mathsf I+\nabla\nabla/k^2)e^{ik|\mathbf r-\mathbf r'|}/4\pi|\mathbf r-\mathbf r'|$:
$\mathbf E(\mathbf r)=k_0^2\int\mathsf G\cdot\Delta\varepsilon\mathbf E\,d^3r'$, an exact identity (Johnson's
Eq.~1). (3) Far field: $\mathsf G\to(\mathsf I-\hat{\mathbf r}\hat{\mathbf r})e^{ikr}e^{-ik\hat{\mathbf r}\cdot\mathbf r'}/4\pi r$.
(4) Adding a feature: subtract the perturbed and unperturbed equations and keep first order; $\delta\mathbf E$ is
defined up to a multiple of $\mathbf E_0$ and the solvability condition fixes $\delta k_0$; numerically one
tracks the resonance. (5) Small feature: the field inside is the quasi static one,
$E_{\rm in}=2\varepsilon_{\rm out}E_{\rm ext}/(\varepsilon_{\rm in}+\varepsilon_{\rm out})$ for a cylinder in a
transverse field and $E_{\rm in}=\varepsilon_{\rm out}E_{\rm out}/\varepsilon_{\rm in}$ for a slab normal to
$\mathbf E$. (6) Projection on a channel uses
$\int\mathbf X^*_{lm}e^{-ik\hat{\mathbf r}\cdot\mathbf r'}d\Omega=4\pi(-i)^lj_l(kr')\mathbf X^*_{lm}(\hat{\mathbf r}')$.
(7) Dipole channel:
$\mathbf p_c/\varepsilon_0=\int\Delta\varepsilon[(j_0-j_2/2)\mathbf E+\tfrac32j_2(\hat{\mathbf r}\cdot\mathbf E)\hat{\mathbf r}]d^3r$,
the exact dipole moment of an extended source. In 2D the same steps give Johnson's Eq.~(3),
$a_m\propto\int J_m(kr)e^{-im\varphi}\Delta\varepsilon E_zd^2r$.

\textbf{Literature} (the full list of v6.2). Johnson, Fan, Mekis, Joannopoulos, APL 78, 3388 (2001): the
mechanism. Karalis, Johnson, Joannopoulos, Opt.\ Lett.\ 29, 2309 (2004): discrete channels, relevant once a
substrate is present. Johnson et al., PRE 65, 066611 (2002): shifting boundaries. Englund, Fushman,
Vu\v ckovi\'c, Opt.\ Express 13, 5961 (2005); Srinivasan, Painter, Opt.\ Express 10, 670 (2002); Quan, Deotare,
Lon\v car, APL 96, 203102 (2010); Quan, Lon\v car, Opt.\ Express 19, 18529 (2011): light cone design of nanobeam
cavities. Nakamura et al., Opt.\ Express 24, 9541 (2016): real space maps of leaky components. Tran, Combri\'e,
De Rossi, PRB 79, 041101 (2009); Portalupi et al., Opt.\ Express 18, 16064 (2010): band folding by a
superimposed period. Rybin et al., PRL 119, 243901 (2017); Bogdanov et al., Adv.\ Photonics 1, 016001 (2019);
Odit et al., Adv.\ Mater.\ 33, 2003804 (2021); Gladyshev, Frizyuk, Bogdanov, PRB 102, 075103 (2020): channel
cancellation in single resonators, symmetry classification. Takata et al., Opt.\ Express 31, 11864 (2023): few
parameter symmetry guided design. Fernandez-Corbaton et al., Opt.\ Express 23, 33044 (2015); Alaee et al.,
Opt.\ Commun.\ 407, 17 (2018): exact multipole moments. Lalanne et al., Laser Photonics Rev.\ 12, 1700113 (2018):
quasinormal mode perturbation theory, the rigorous form of the first order split. Loulas et al., ACS Photonics
12, 1524 (2025): 2D vector multipoles. Chao et al., Opt.\ Express 34, 7337 (2026): bounds based initialisation of
inverse design. La et al., Laser Photonics Rev.\ (2026), doi 10.1002/lpor.202501030: light ultrasound interaction
in a $\pi$ shifted fibre grating, the sensor side. Sci.\ Rep.\ (2026), doi 10.1038/s41598-026-51041-9: SiN Bragg
mirrors, apodisation, coupling versus mirror loss. Ansys, ``Far field projections in FDTD overview''. No
publication was found applying channel cancellation to a sidewall Bragg cavity at fixed mode width.
\end{document}
